
\section{\BookRecordIndex{T0142}格罗莫夫拓扑}\label{sec:D-5}

本节内在地描述\BookRecordIndex{T0246}模空间 $\overline{\mathcal M}_{0,n}$ 上的拓扑，并证明它与由嵌入到二维球面乘积所定义的外在拓扑等价。以下定义是\BookRecordIndex{T0340}稳定映射的\BookRecordIndex{T0141}格罗莫夫收敛在目标流形为单点时的特例（见定义~\ref{def:5-5-1}）。

\begin{definition}\label{def:D-5-1}
设
\[
  \mathbf z^\nu=\bigl(\{z^\nu_{\alpha\beta}\}_{\alpha E^\nu\beta},
  \{\alpha_i^\nu,z_i^\nu\}_{1\le i\le n}\bigr)\in SC_{0,n}
\]
是一列\BookRecordIndex{T0339}稳定曲线，并设
\[
  \mathbf z=\bigl(\{z_{\alpha\beta}\}_{\alpha E\beta},
  \{\alpha_i,z_i\}_{1\le i\le n}\bigr).
\]
若对充分大的 $\nu$，存在满射的\BookRecordIndex{T0305}树同态 $f^\nu:T\to T^\nu$ 及一族 \BookRecordIndex{T0040}Möbius 变换 $\{\phi_\alpha^\nu\}_{\alpha\in T}^{\nu\in\mathbb N}$，满足以下条件，就称 $\mathbf z^\nu$ \emph{\BookRecordIndex{T0141}格罗莫夫收敛}于 $\mathbf z$。

\noindent\textup{（重标度）}若 $\alpha,\beta\in T$ 且 $\alpha E\beta$，并且 $\nu_j$ 是满足 $f^{\nu_j}(\alpha)=f^{\nu_j}(\beta)$ 的子序列，则
\[
  \phi_{\alpha\beta}^{\nu_j}:=(\phi_\alpha^{\nu_j})^{-1}\circ\phi_\beta^{\nu_j}
  \longrightarrow z_{\alpha\beta}
\]
在 $S^2\setminus\{z_{\beta\alpha}\}$ 的紧子集上一致收敛。

\noindent\textup{（\BookRecordIndex{T0191}节点）}若 $\alpha,\beta\in T$ 且 $\alpha E\beta$，并且 $\nu_j$ 是满足 $f^{\nu_j}(\alpha)\ne f^{\nu_j}(\beta)$ 的子序列，则
\[
  z_{\alpha\beta}=\lim_{j\to\infty}
  (\phi_\alpha^{\nu_j})^{-1}
  \bigl(z_{f^{\nu_j}(\alpha)f^{\nu_j}(\beta)}^{\nu_j}\bigr).
\]

\noindent\textup{（\BookRecordIndex{T0066}标记点）}对每个 $i$，有 $\alpha_i^\nu=f^\nu(\alpha_i)$，并且
\[
  z_i=\lim_{\nu\to\infty}(\phi_{\alpha_i}^\nu)^{-1}(z_i^\nu).
\]
\end{definition}

因为每棵树 $T$ 的满射像只有有限多个同构类，任何\BookRecordIndex{T0141}格罗莫夫收敛序列都是有限个子序列的并；在每个子序列上，树 $T^\nu$ 和映射 $f^\nu$ 均同构。上面（重标度）与（\BookRecordIndex{T0191}节点）公理涉及的正是这些子序列。

本节的主要结论如下。

\begin{theorem}\label{thm:D-5-2}
$\overline{\mathcal M}_{0,n}$ 上存在唯一的拓扑，使其收敛序列恰为\BookRecordIndex{T0141}格罗莫夫收敛序列。此拓扑与定理~\ref{thm:D-4-2} 的嵌入所诱导的 $\overline{\mathcal M}_{0,n}$ 上的拓扑一致。
\end{theorem}

在此拓扑下，$\overline{\mathcal M}_{0,n}$ 与 $\overline M_{0,n}$ 同胚，因而具有紧致度量空间的结构，实际上还是紧致光滑代数簇。

\begin{proposition}\label{prop:D-5-3}
设 $\mathbf z^\nu,\mathbf z\in SC_{0,n}$。以下条件等价：
\begin{enumerate}[label=(\roman*)]
\item $\mathbf z^\nu$ \BookRecordIndex{T0141}格罗莫夫收敛于 $\mathbf z$；
\item 对充分大的 $\nu$，存在满射\BookRecordIndex{T0305}树同态 $f^\nu:T\to T^\nu$ 和 \BookRecordIndex{T0040}Möbius 变换 $\phi_\alpha^\nu\in G$，使 $f^\nu(\alpha_i)=\alpha_i^\nu$，并且对 $i\in\{1,\ldots,n\}$、$\alpha\in T$，
\begin{equation}\label{eq:D-5-1}
  z_{\alpha i}=\lim_{\nu\to\infty}(\phi_\alpha^\nu)^{-1}
  \bigl(z_{f^\nu(\alpha)i}^\nu\bigr);
\end{equation}
\item 对任意四个互异的整数 $i,j,k,\ell\in\{1,\ldots,n\}$，
\begin{equation}\label{eq:D-5-2}
  w_{ijk\ell}(\mathbf z)=\lim_{\nu\to\infty}w_{ijk\ell}(\mathbf z^\nu).
\end{equation}
\end{enumerate}
\end{proposition}

\begin{proof}[命题~\ref{prop:D-5-3} \BookRecordIndex{T0326}推出定理~\ref{thm:D-5-2}]
由引理~\ref{lem:5-6-4}，为了证明 $\overline{\mathcal M}_{0,n}$ 上存在唯一拓扑，其收敛序列恰为\BookRecordIndex{T0141}格罗莫夫收敛序列，只须核验（常值）、（子序列）、（子子序列）、（对角线）及（唯一性）这五条公理。命题~\ref{prop:D-5-3} 的 (i) 与 (iii) 等价，说明序列\BookRecordIndex{T0141}格罗莫夫收敛当且仅当其在 $\overline M_{0,n}$ 中的像收敛，故这些公理成立。所得拓扑与拉回拓扑具有相同的收敛序列，因而必定一致。
\end{proof}

\begin{proof}[命题~\ref{prop:D-5-3} 的证明]
先证 (i) 蕴含 (ii)。设 $\mathbf z^\nu$ \BookRecordIndex{T0141}格罗莫夫收敛于 $\mathbf z$，反设某个 $\alpha\in T$ 与 $i\in\{1,\ldots,n\}$ 不满足~\eqref{eq:D-5-1}。取子序列后可设
\begin{equation}\label{eq:D-5-3}
  \lim_{\nu\to\infty}(\phi_\alpha^\nu)^{-1}
  \bigl(z_{f^\nu(\alpha)i}^\nu\bigr)=z_0\ne z_{\alpha i}.
\end{equation}
再取子序列，可设 $T^\nu=T'$ 且\BookRecordIndex{T0305}树同态 $f^\nu=f:T\to T'$ 均不依赖于 $\nu$。写 $f(\alpha)=:\alpha'$；这与标记的记号\BookRecordIndex{T0348}相容，因为 $f(\alpha_i)=\alpha_i'$，其中 $\alpha_i\in T$、$\alpha_i'\in T'$ 分别是承载第 $i$ 个\BookRecordIndex{T0066}标记点的顶点。

由~\eqref{eq:D-5-3} 和（\BookRecordIndex{T0066}标记点）公理，$\alpha\ne\alpha_i$。取从 $\gamma_{-\ell}=\alpha_i$ 到 $\gamma_k=\alpha$ 的链
\[
  \gamma_{-\ell},\ldots,\gamma_{-1},\gamma_0,\ldots,\gamma_k,
\]
使标号满足 $\gamma_j'=\alpha'$（$j\ge0$）及 $\gamma_j'\ne\alpha'$（$j<0$）。首先断言
\begin{equation}\label{eq:D-5-4}
  z_{\gamma_0i}=\lim_{\nu\to\infty}(\phi_{\gamma_0}^\nu)^{-1}
  (z_{\gamma_0'i}^\nu).
\end{equation}
若 $\alpha'=\alpha_i'$，则 $\gamma_0=\alpha_i$，故 $z_{\gamma_0i}=z_i$ 且 $z_{\gamma_0'i}^\nu=z_i^\nu$；此时~\eqref{eq:D-5-4} 来自定义~\ref{def:D-5-1} 的（\BookRecordIndex{T0066}标记点）公理。若 $\alpha'\ne\alpha_i'$，则 $z_{\gamma_0i}=z_{\gamma_0\gamma_{-1}}$，且 $z_{\gamma_0'i}^\nu=z_{\gamma_0'\gamma_{-1}'}^\nu$；此时~\eqref{eq:D-5-4} 来自（\BookRecordIndex{T0191}节点）公理。

有了~\eqref{eq:D-5-4}，我们归纳证明
\begin{equation}\label{eq:D-5-5}
  z_{\gamma_ji}=\lim_{\nu\to\infty}(\phi_{\gamma_j}^\nu)^{-1}
  (z_{\gamma_j'i}^\nu),\qquad j=1,\ldots,k.
\end{equation}
设它对 $j\le k-1$ 成立。注意 $z_{\gamma_ji}=z_{\gamma_j\gamma_{j-1}}\ne z_{\gamma_j\gamma_{j+1}}$。由（重标度）公理，$(\phi_{\gamma_{j+1}}^\nu)^{-1}\circ\phi_{\gamma_j}^\nu$ 在 $S^2\setminus\{z_{\gamma_j\gamma_{j+1}}\}$ 的紧子集上一致收敛于 $z_{\gamma_{j+1}\gamma_j}=z_{\gamma_{j+1}i}$。因此得到 $j+1$ 时的~\eqref{eq:D-5-5}，归纳完成。取 $j=k$ 即与~\eqref{eq:D-5-3} 矛盾，故 (i) 蕴含 (ii)。

现证 (ii) 蕴含 (iii)。仍只须考虑对每个 $\nu$ 都有 $T^\nu=T'$、$f^\nu=f:T\to T'$ 的情形。任取互异的 $i,j,k,\ell\in\{1,\ldots,n\}$，选择 $\alpha\in T$ 使 $(z_{\alpha i},z_{\alpha j},z_{\alpha k},z_{\alpha\ell})\notin\Delta_3$。由假设，$\alpha'=f(\alpha)$ 时，$z_{\alpha i}=\lim_{\nu\to\infty}(\phi_\alpha^\nu)^{-1}(z_{\alpha'i}^\nu)$；其余三个指标同理。因此对充分大的 $\nu$，$(z_{\alpha'i}^\nu,z_{\alpha'j}^\nu,z_{\alpha'k}^\nu,z_{\alpha'\ell}^\nu)\notin\Delta_3$，且
\begin{align*}
  w_{ijk\ell}(\mathbf z)
  &=w(z_{\alpha i},z_{\alpha j},z_{\alpha k},z_{\alpha\ell})\\
  &=\lim_{\nu\to\infty}w\bigl((\phi_\alpha^\nu)^{-1}(z_{\alpha'i}^\nu),\ldots,
      (\phi_\alpha^\nu)^{-1}(z_{\alpha'\ell}^\nu)\bigr)\\
  &=\lim_{\nu\to\infty}w(z_{\alpha'i}^\nu,z_{\alpha'j}^\nu,
      z_{\alpha'k}^\nu,z_{\alpha'\ell}^\nu)\\
  &=\lim_{\nu\to\infty}w_{ijk\ell}(\mathbf z^\nu).
\end{align*}
这证明了~\eqref{eq:D-5-2}。

再证 (iii) 蕴含 (ii)。若 (iii) 成立，则充分大的 $\nu$ 满足
\[
  w_{ijk\ell}(\mathbf z)\ne\infty\quad\Longrightarrow\quad
  w_{ijk\ell}(\mathbf z^\nu)\ne\infty,
\]
从而 $S(\mathbf z^\nu)\subset S(\mathbf z)$。由定理~\ref{thm:D-2-6}(iv)，对这样的 $\nu$，存在唯一的满射\BookRecordIndex{T0305}树同态 $f^\nu:T\to T^\nu$，使所有 $i$ 均满足 $f^\nu(\alpha_i)=\alpha_i^\nu$。下文记 $\alpha^\nu:=f^\nu(\alpha)\in T^\nu$。对每个 $\alpha\in T$ 选取 $i,j,k\in\{1,\ldots,n\}$，使 $\alpha=\alpha(i,j,k;T)$。由定理~\ref{thm:D-2-6}(iv)，$\alpha^\nu=\alpha(i,j,k;T^\nu)$。依注~\ref{rem:D-3-3}，这意味着 $z_{\alpha i},z_{\alpha j},z_{\alpha k}$ 两两不同，且对充分大的 $\nu$，$z_{\alpha^\nu i}^\nu,z_{\alpha^\nu j}^\nu,z_{\alpha^\nu k}^\nu$ 也两两不同。因而存在唯一的 \BookRecordIndex{T0040}Möbius 变换 $\phi_\alpha^\nu\in G$，满足
\[
  \phi_\alpha^\nu(z_{\alpha i})=z_{\alpha^\nu i}^\nu,\qquad
  \phi_\alpha^\nu(z_{\alpha j})=z_{\alpha^\nu j}^\nu,\qquad
  \phi_\alpha^\nu(z_{\alpha k})=z_{\alpha^\nu k}^\nu.
\]
断言这一序列满足~\eqref{eq:D-5-1}。确实，
\begin{align*}
  w_{ijk\ell}(\mathbf z^\nu)
  &=w(z_{\alpha^\nu i}^\nu,z_{\alpha^\nu j}^\nu,
      z_{\alpha^\nu k}^\nu,z_{\alpha^\nu\ell}^\nu)\\
  &=\frac{(z_{\alpha j}-z_{\alpha k})(\zeta_{\alpha\ell}^\nu-z_{\alpha i})}
          {(z_{\alpha i}-z_{\alpha j})(z_{\alpha k}-\zeta_{\alpha\ell}^\nu)},
  \qquad
  \zeta_{\alpha\ell}^\nu:=(\phi_\alpha^\nu)^{-1}
       (z_{\alpha^\nu\ell}^\nu).
\end{align*}
\sourcefootnote{原书实体页 657 此分式的分母第二因子印为 $z_{\alpha i}-\zeta_{\alpha\ell}^\nu$。由前一行的\BookRecordIndex{T0185}交比 $w(z_{\alpha i},z_{\alpha j},z_{\alpha k},\zeta_{\alpha\ell}^\nu)$ 及 D.1 节的\BookRecordIndex{T0185}交比定义，该因子应为 $z_{\alpha k}-\zeta_{\alpha\ell}^\nu$，此处据此补正。}
利用 $w_{ijk\ell}(\mathbf z^\nu)\to w_{ijk\ell}(\mathbf z)$，可得
\[
  z_{\alpha\ell}=\lim_{\nu\to\infty}\zeta_{\alpha\ell}^\nu
\]
对所有 $\alpha\in T$、$\ell\in\{1,\ldots,n\}$ 成立，故 (ii) 得证。

最后证 (ii) 蕴含 (i)。若 (ii) 成立，定义~\ref{def:D-5-1} 的（\BookRecordIndex{T0066}标记点）公理显然成立（在~\eqref{eq:D-5-1} 中取 $\alpha=\alpha_i$）。为验证（\BookRecordIndex{T0191}节点）公理，选 $i$ 使 $z_{\alpha\beta}=z_{\alpha i}$，于是 $z_{\alpha^\nu\beta^\nu}^\nu=z_{\alpha^\nu i}^\nu$，再应用~\eqref{eq:D-5-1}。为验证（重标度）公理，设 $\alpha E\beta$ 且 $f^\nu(\alpha)=f^\nu(\beta)=\alpha^\nu$ 对所有 $\nu$ 成立。选指标 $j,j'$，使 $z_{\beta\alpha},z_{\beta j},z_{\beta j'}$ 两两不同，因而 $z_{\alpha j}=z_{\alpha j'}=z_{\alpha\beta}$。由~\eqref{eq:D-5-1}，
\[
  \lim_{\nu\to\infty}(\phi_\beta^\nu)^{-1}(z_{\alpha^\nu j}^\nu)
   =z_{\beta j}\ne z_{\beta j'}
   =\lim_{\nu\to\infty}(\phi_\beta^\nu)^{-1}(z_{\alpha^\nu j'}^\nu),
\]
\[
  \lim_{\nu\to\infty}(\phi_\alpha^\nu)^{-1}(z_{\alpha^\nu j}^\nu)
   =z_{\alpha\beta}
   =\lim_{\nu\to\infty}(\phi_\alpha^\nu)^{-1}
       (z_{\alpha^\nu j'}^\nu).
\]
交换 $\alpha,\beta$ 的角色，还可找到指标 $i$，使
\[
  \lim_{\nu\to\infty}(\phi_\alpha^\nu)^{-1}(z_{\alpha^\nu i}^\nu)
    =z_{\alpha i}\ne z_{\alpha\beta},\qquad
  \lim_{\nu\to\infty}(\phi_\beta^\nu)^{-1}(z_{\alpha^\nu i}^\nu)
    =z_{\beta\alpha}.
\]
由引理~\ref{lem:D-1-4}，$(\phi_\alpha^\nu)^{-1}\circ\phi_\beta^\nu$ 在 $S^2\setminus\{z_{\beta\alpha}\}$ 的紧子集上一致收敛于 $z_{\alpha\beta}$。命题得证。
\end{proof}
